\section{Detailed prior-art comparison}
\label{app:lineage}

Kempthorne's Section 3 begins with a finite prior, locally maximizes Bayes risk
over weights and locations, globally maximizes the Bayes-rule risk, stops when
maximum risk equals Bayes risk, and otherwise inserts the maximizer
\cite[pp.~175--177]{kempthorne1987}.  He also notes that zero weights or
coincident locations can reduce support.  His implementation brackets local
risk maxima by a parameter grid and refines them by golden-section search
\cite[p.~177]{kempthorne1987}.  The present verifier adds outward rounding
and complete continuous branch coverage to that optimization lineage.

Gourdin et al.\ terminate their Lipschitz global optimizer when the gap
between an upper bound and incumbent falls below a chosen tolerance
\cite[pp.~19--21]{gourdin1994}.  Their bounded-normal discussion fixes a
finite-dimensional prior form justified for a small range of $m$
\cite[pp.~25--26]{gourdin1994}.  The present solver combines adaptive support
generation with an interval-certified continuous risk oracle.

Dytso et al.\ \cite{dytso2022} give support-cardinality reductions under a
finite observation alphabet and use projected gradient ascent in numerical
examples.  Their stated finite-alphabet theorem addresses a different
observation space from the scalar Gaussian model.  Kent et al.\
\cite{kent2021} assume derivative and local linear-minimization oracles and
use stochastic subroutines.  Guggenberger and Huang
\cite{guggenberger2025} obtain iterationwise brackets in a
finite/discretized fictitious-play problem with grid-based best responses.
These results locate the current work within broader finite-mixture and
minimax-optimization methods.

Maruyama and Matsuda \cite{maruyama2025} provide the closest located
interval-assisted normal-mean minimaxity precedent.  They verify inequalities
for a Bayes estimator under a fixed half-Cauchy prior in a multivariate
problem.  The certified object here is instead the Bayes/maximum-risk gap for
a least-favorable prior in the bounded scalar problem.

This comparison is source-grounded and records the closest antecedents used in
the paper.  Future revisions should incorporate any closer continuous,
outward-rounded bounded-normal certificate brought to the author's attention.
